\section{Experimental Setup}
\label{sec:setup}

\paragraph{LLMs.} Our data submissions use two judge models. The first is
Qwen3-30B-A3B-Instruct-2507 from the Qwen3 family~\cite{yang2025qwen3}, with its
weights in FP8, a low-precision floating-point format. We call it Qwen3-30B. It is a mixture-of-experts
model. Each token passes through only a few expert sub-networks, so about 3B
parameters are active per token. This 2507 Instruct release has no thinking mode,
so it writes no reasoning block. The second, gpt-oss-20b~\cite{openai2025gptoss},
is a reasoning model. We also ran Qwen2.5-7B-Instruct~\cite{qwen2025qwen25} on one
pilot, because the organizers named it among the models they would use. We served
all three with vLLM 0.13.0~\cite{kwon2023vllm}, an open-source LLM serving engine,
on one NVIDIA RTX 6000 Ada GPU with 48\,GB of memory. We allowed a context of
32,768 tokens and up to 64 concurrent sequences. A full run with Qwen3-30B
took 30 minutes on rag26 and 14 minutes on ragtime26. gpt-oss-20b is slower. It
took 99 minutes on ragtime25, at about 40 seconds per call.

\paragraph{Reply budgets.} Our main data submission is \run{nuggetrecall-best} (Table~\ref{tab:submissions}). In it and on the pilots, our Qwen3-30B and Qwen2.5-7B runs capped replies at
1,536 tokens for the key and 1,024 for judging. We later raised both to 4,096 for
reasoning models. Neither Qwen model reasons before it answers. vLLM also recorded
no Qwen3-30B reply that reached the old limit. Our gpt-oss-20b runs used 4,096
tokens on every attempt, except the rag26 run, which had the growing budget of
Section~\ref{sec:budget}.

\paragraph{Runs by the organizers.} The organizers ran our code submissions on
TIRA with models they chose. We made two \judge{NuggetRecall} code submissions, \run{rank-dioptre} and \run{grave-node} (Section~\ref{sec:submissions}). On the pilots the organizers ran \run{rank-dioptre}, which keeps
4,096 tokens on every attempt. On rag25-gen they paired it with gpt-oss-120b. On
dragun-repgen they ran it with Qwen2.5-7B-Instruct, gpt-oss-20b and gpt-oss-120b.
Our \judge{Corroborate} submission, \run{lead-racer}, ran on CPU on all three
pilots. On the test sets they ran \run{lead-racer}, plus \run{rank-dioptre} and
\run{grave-node} with the same three models. \run{grave-node} also got a fourth
model, listed on TIRA as Qwen3.8-27B-FP8. Tables mark organizer runs with \org.

\paragraph{Encoders.} \judge{Corroborate} uses bge-base-en-v1.5 to embed
sentences and the nli-deberta-v3-base cross-encoder to check entailment. Both run
on a GPU when one has at least 6\,GB free, or else on CPU with eight threads.

\paragraph{Validation.} We fitted nothing on the test sets. Every choice we made by
score came from the pilots. Three \judge{Corroborate} constants have other
sources. The similarity threshold $\theta$ and the partial-credit cutoff come from the synthetic
dataset (Section~\ref{sec:corroborate}). The specificity floor is based on the
encoder's range (Section~\ref{sec:floor}). One engineering change followed an
observation on test data. On the first pass, 27\% of the rag26 judgments by gpt-oss-20b came back
unreadable. So we added the growing reply budget of Section~\ref{sec:budget} and
asked the organizers to treat the code submission carrying it as our main one. We
saw only that count and vLLM's count of replies cut off at the length limit. We
never saw a response, a grade or any truth data.

\judge{NuggetRecall} is not deterministic across runs (Section~\ref{sec:noise}).
\judge{Corroborate} is. Its output on our GPU matched the organizers' CPU run to
within 0.002. Section~\ref{sec:noise} measures how much \judge{NuggetRecall} varies from run to run.

\paragraph{Code submission.} Our machines have no container runtime such as
Docker, so we build and submit the containers from GitHub Actions. Before each
submission the workflow runs the test suite. Next it fills the prompt
cache by running the judge once against a stub LLM endpoint, a stand-in for a
real model. Then it runs the container with the endpoint switched off. To pass, the container must
reproduce its own leaderboard from that cache. That leaderboard must also actually
rank systems apart. The starter base image ships Python 3.11.0rc1, a release
candidate. On it, a current transformers and PyTorch stack fails at import, so our
image installs its own Python 3.12.
